\section{API et bases de données}
\label{sec:backend}

Cette section décrit les choix techniques retenus pour le backend de l'application : l'architecture mise en place, la structure de la base de données et le fonctionnement de l'API.

\subsection{Architecture générale du backend}

Le backend est développé avec \textbf{FastAPI}, un framework Python moderne taillé pour la création d'API REST. Pour la persistance des données, nous avons choisi \textbf{PostgreSQL}, piloté via l'ORM \textbf{SQLAlchemy} qui nous permet de manipuler la base de données sous forme d'objets Python, ce qui simplifie considérablement la lecture et la maintenance du code. L'ensemble de l'infrastructure est containerisé avec \textbf{Docker}, ce qui nous garantit un environnement identique entre le développement et la production.

La structure du projet suit une architecture modulaire en couches :
\begin{itemize}
    \item \texttt{models/} : Définition des entités et des tables via SQLAlchemy.
    \item \texttt{schemas/} : Validation et typage des données entrantes et sortantes via Pydantic.
    \item \texttt{routers/} : Définition des points d'accès (endpoints) de l'API.
    \item \texttt{utils/} : Fonctions transversales partagées, notamment la sécurité et le hachage.
\end{itemize}

\subsection{Base de données}

La base de données est le socle du système : elle centralise les données des utilisateurs et assure la persistance de leur activité de navigation.

\subsubsection{Modélisation des données}

Nous avons modélisé les entités autour des besoins concrets de l'application : les comptes utilisateurs, leurs vélos, les itinéraires calculés, l'historique de navigation et les signalements communautaires. Chaque entité correspond à une table PostgreSQL, créée automatiquement au démarrage du serveur via SQLAlchemy.

\subsubsection{Schéma relationnel}

Les entités sont liées entre elles par des relations qui reflètent la logique métier de l'application :
\begin{itemize}
    \item Un \textbf{utilisateur} peut posséder plusieurs \textbf{vélos} (relation 1-N).
    \item Un \textbf{utilisateur} peut sauvegarder plusieurs \textbf{itinéraires} (relation 1-N).
    \item Un \textbf{itinéraire} peut être associé à plusieurs entrées dans l'\textbf{historique} de navigation (relation 1-N).
    \item Un \textbf{utilisateur} peut émettre plusieurs \textbf{signalements} de dangers (relation 1-N).
\end{itemize}

\subsubsection{Description des tables}

Voici le rôle de chaque table dans la base de données :

\begin{itemize}
    \item \textbf{Table \texttt{users}} : Stocke les données de profil, dont les identifiants de connexion hachés et les informations personnelles.
    \item \textbf{Table \texttt{bikes}} : Contient les caractéristiques techniques des vélos (type, assistance électrique) rattachés à chaque utilisateur.
    \item \textbf{Table \texttt{routes}} : Enregistre les trajets calculés avec leurs coordonnées, distance, durée estimée et indice de sécurité.
    \item \textbf{Table \texttt{user\_history}} : Retrace l'activité de navigation en associant chaque trajet effectué à un horodatage.
    \item \textbf{Table \texttt{reports}} : Conserve les signalements géolocalisés émis par la communauté (accidents, travaux, obstacles).
    \item \textbf{Tables \texttt{badges} et \texttt{user\_badges}} : Posent les bases d'un système de gamification prévu pour une version ultérieure.
\end{itemize}

\subsection{Gestion du trafic en temps réel}

L'un des aspects critiques de la sécurité cycliste est l'évitement des zones de forte congestion automobile. Un axe embouteillé n'est pas seulement synonyme de lenteur : il concentre stress, manœuvres imprévisibles et risques d'accident pour le cycliste. Notre système intègre donc une couche de données trafic en temps réel pour adapter les informations affichées aux conditions de circulation actuelles.

\subsubsection{Source de données : l'API AVATAR du Cerema}

Pour alimenter cette fonctionnalité, nous utilisons l'\textbf{API AVATAR du Cerema}, qui recense les points de comptage fixes déployés sur la métropole bordelaise — boucles électromagnétiques et radars. Pour chaque capteur, trois indicateurs sont disponibles : la vitesse moyenne ($v$, en km/h), le débit horaire ($q$, en véhicules/heure) et le taux d'occupation ($t$, en \%). Ces données sont récupérées par le backend et exposées aux clients via l'endpoint \texttt{/traffic/}.

\subsubsection{Stratégie de mise à jour et gestion du cache}

La liste des capteurs Cerema est géographiquement stable : les points de comptage ne changent pas d'emplacement d'une heure à l'autre. Nous la mettons donc en cache en mémoire avec une durée de vie (\textit{TTL}) d'une heure, ce qui évite de la récupérer inutilement à chaque requête. En revanche, les mesures de trafic elles-mêmes sont interrogées dynamiquement à chaque appel, en se limitant aux données des deux dernières heures. Cette fenêtre glissante garantit que les niveaux affichés reflètent la circulation du moment, sans remonter à des mesures obsolètes.

\subsubsection{Calcul de l'indice de danger}

À partir des mesures brutes de chaque capteur, nous calculons un score de danger $D$ selon la formule pondérée suivante :
\begin{equation}
    D = s_v \times 0{,}6 + s_q \times 0{,}4
\end{equation}
$s_v$ est le score de vitesse et $s_q$ le score de débit, tous deux normalisés entre 0 (aucun danger) et 1 (danger maximal). Les seuils sont les suivants : $s_v = 0$ si $v \geq 70$\,km/h et $s_v = 1$ si $v < 30$\,km/h ; $s_q = 0$ si $q < 200$\,véh./h et $s_q = 1$ si $q \geq 800$\,véh./h.

Nous avons choisi de donner plus de poids à la vitesse (60\,\%) qu'au débit (40\,\%) pour une raison simple : du point de vue d'un cycliste, ce qui rend un axe dangereux, c'est avant tout la vitesse à laquelle les véhicules circulent. Un trafic dense mais lent reste gérable ; en revanche, des voitures rapides sur une voie sans aménagement cyclable exposent le cycliste à des risques bien plus sérieux. Le débit reste un indicateur utile pour détecter la saturation de la chaussée, mais il est secondaire face à l'effet de la vitesse. Le score final classe chaque point de comptage en trois niveaux : \textit{fluide} ($D < 0{,}35$, vert), \textit{dense} ($D < 0{,}65$, orange) ou \textit{critique} ($D \geq 0{,}65$, rouge). Faute de données suffisantes, le niveau reste \textit{indéterminé} (gris).

\subsection{API REST}

L'API constitue l'interface entre les clients (web et mobile) et la logique serveur. Elle est construite autour de FastAPI et expose un ensemble d'endpoints organisés par ressource.

\subsubsection{Technologies utilisées}

FastAPI génère automatiquement une documentation interactive via Swagger UI, ce qui a facilité nos tests pendant le développement. La validation des données entrantes est assurée par \textbf{Pydantic} : chaque requête est vérifiée avant d'atteindre la logique métier, ce qui réduit les risques d'erreurs en cascade.

Pour la sécurité des échanges, nous nous appuyons sur deux standards : \textbf{JWT} (JSON Web Token) pour la gestion des sessions, et \textbf{Argon2} pour le hachage des mots de passe, un algorithme reconnu pour sa résistance aux attaques par force brute.

\subsubsection{Authentification et sécurité}

L'authentification repose sur le standard \textbf{OAuth2 avec Bearer Token}. À la connexion, le serveur délivre un jeton JWT que le client doit inclure dans l'en-tête de chaque requête protégée.

Nous distinguons deux niveaux d'accès :
\begin{itemize}
    \item \textbf{Accès authentifié} : obligatoire pour les actions sensibles — modification du profil, gestion des vélos, consultation de l'historique et émission de signalements.
    \item \textbf{Accès public} : le calcul d'itinéraire est accessible sans connexion, mais la sauvegarde automatique du trajet ne s'active que pour les utilisateurs connectés.
\end{itemize}

\subsubsection{Endpoints principaux}

L'API expose les endpoints suivants, organisés par ressource :

\begin{itemize}
    \item \textbf{\texttt{/users}} : Gestion du cycle de vie des comptes, de l'inscription à la mise à jour des informations personnelles.
    \item \textbf{\texttt{/bikes}} : Administration des vélos de l'utilisateur, avec un contrôle strict garantissant qu'un utilisateur ne peut accéder qu'à ses propres équipements.
    \item \textbf{\texttt{/routes}} : Calcul d'itinéraires à partir de coordonnées géographiques et des préférences de l'utilisateur (type de vélo, niveau de sécurité souhaité).
    \item \textbf{\texttt{/history}} : Consultation des trajets passés de l'utilisateur connecté.
    \item \textbf{\texttt{/stats}} : Synthèse de l'activité de l'utilisateur — distances parcourues, durée cumulée, répartition entre trajets rapides et sécurisés.
    \item \textbf{\texttt{/reports}} : Remontée de signalements terrain par les usagers. La lecture est publique pour favoriser la sécurité collective ; l'écriture requiert une authentification.
\end{itemize}
